% !TEX root = ../main.tex
% !TEX program = lualatex
\documentclass[../main.tex]{subfiles}
\IfSubfilesClassLoaded{\externaldocument{\subfix{../build/main}}}{}
% =============================================================================
% CHAPTER 1: COMMAND CHAINS AND GOVERNANCE
% DESCRIPTION: Operational vs. acquisition chains, Title 10 distinctions, and governance.
% =============================================================================

\begin{document}
\IfSubfilesClassLoaded{\chapter{EDO Study Guide}}{}
%====================
% Contents here
%====================
% === COMMAND CHAINS & GOVERNANCE ===
\section{Command Chains and Governance}

\subsection{Foundations and Goldwater-Nichols}
\begin{description}
  \item[\textbf{Constitutional basis.}] The Constitution establishes three branches of government: Article I (Congress) authorizes raising and supporting armies, providing and maintaining a navy, and declaring war; Article II vests civilian control in the President as Commander in Chief; Article III creates the judiciary~\autocite{edo-2-1-1-command-structures-2025}.
  \item[\textbf{Pre-1986 context.}] Service leaders historically carried both force-provider and force-commander responsibilities. The National Security Act of 1947 formalized the Joint Chiefs and created the National Military Establishment, and the 1949 amendment renamed it the Department of War~\autocite{edo-2-1-1-command-structures-2025}.
  \item[\textbf{Goldwater-Nichols (1986).}] The Act strengthened the \ac{cjcs} as the principal military advisor, increased \ac{ccdr} authority, and institutionalized jointness in operations, education, and budgeting~\autocite{Goldwater-Nichols-1986,edo-2-1-1-command-structures-2025}.
\end{description}

\subsection{Operational vs.\ Administrative Chains}
\begin{description}
  \item[\textbf{Operational chain (OPCON).}] The operational chain runs President \(\to\) \ac{secdef} \(\to\) \acp{ccdr}; \ac{opcon} authorities flow through that chain, and Service chiefs are not in the operational chain of command~\autocite{USC-113,USC-162}.
  \item[\textbf{Administrative chain (ADCON).}] \ac{adcon} aligns readiness and resourcing: Service Secretaries and chiefs \emph{man, train, and equip} forces; for the Navy this runs through \ac{secnav} and \ac{cno} down to \acp{syscom} and operating forces~\autocite{USC-8032,edo-2-1-1-command-structures-2025}.
\end{description}

Figure \ref{fig:chain_of_command_ops} shows the Title 10 chain of command. Figure \ref{fig:chain_of_command_org_chart} shows the Navy chain from the President to Echelon III commands.

\begin{figure}[H]
  \centering\color{SteelDeep}
  \includegraphics[width=\linewidth]{Command Structure.jpeg}
  \caption[Title 10 Chain of Command]{Title 10 Chain of Command.~\srcCite{edo-2-1-1-command-structures-2025}}
  \label{fig:chain_of_command_ops}
\end{figure}

\begin{figure}[H]
  \centering\color{SteelDeep}
  \includegraphics[width=\linewidth]{Org_Chart.pdf}
  \caption[Organization command structure]{Organization command structure from the President through Echelon III.~\srcCite{edo-2-1-1-command-structures-2025}}
  \label{fig:chain_of_command_org_chart}
\end{figure}

\subsection{DoW and DoN Components}
\begin{description}
  \item[\textbf{DoW components.}] Core elements include the \ac{osd}, Joint Staff, Military Departments, Unified \acp{ccmd}, Defense Agencies (e.g., \ac{mda}), and \ac{dod} Field Activities~\autocite{edo-2-1-1-command-structures-2025}.
  \item[\textbf{DoN components.}] The \ac{don} includes the Secretariat (\ac{secnav} and principal advisors), \ac{cno}/\ac{opnav} staff, the Shore Establishment (\acp{syscom} and field activities), and the Operating Forces~\autocite{edo-2-1-1-command-structures-2025}.
\end{description}

\subsection{Strategic Direction (NSS and NDS)}
The \ac{nss} sets national-level priorities; the \ac{nds} translates them into defense priorities, posture, and force-development guidance. The 2022 \ac{nds} emphasizes defending the homeland, deterring strategic attacks, deterring aggression while prepared to prevail, and building a resilient joint force and defense ecosystem; it advances these priorities through integrated deterrence, campaigning, and building enduring advantage. Current as of 2026-01-11~\autocite{NSS-2022,NDS-2022}.

\subsection{SECNAV Principal Advisors}
\begin{longtblr}[%
    caption = {SECNAV principal advisors},
    label   = {tab:secnav-principal-advisors},
    entry   = {SECNAV principal advisors},
    remark{\tSource} = {\tabCite{edo-2-1-1-command-structures-2025}}
  ]{colspec = {@{} X[0.6] X[1.4] @{}}}
  \toprule
  {Office} & {Focus and responsibilities} \\
  \midrule
  Under Secretary of the Navy & Leads the Secretariat and serves as the principal advisor to the Secretary of the Navy for all Department of the Navy matters. \\
  Assistant Secretary of the Navy (Manpower and Reserve Affairs) & Manpower policy, personnel readiness, quality of life, and workforce programs across the Navy and Marine Corps. \\
  Assistant Secretary of the Navy (Installations and Environment) & Real property, housing, facilities, environmental protection, and BRAC processes. \\
  Assistant Secretary of the Navy (Financial Management and Comptroller) & Budget formulation and execution; principal financial advisor and comptroller for the Department of the Navy. \\
  Assistant Secretary of the Navy (Research, Development and Acquisition) & Acquisition policy, program oversight, and life-cycle management for Navy and Marine Corps systems. \\
  \bottomrule
\end{longtblr}

The \ac{asnrdanda} is the \ac{don} \ac{cae} and the \ac{nae}, with authority and accountability for acquisition policy and program oversight~\autocite{SECNAV5400-15D}.

\subsection{Acquisition Governance (Department of the Navy)}

\begin{description}
  \item[\textbf{Policy authority.}] \ac{asnrdanda} is the \ac{sae} for the \ac{don} and sets policy for \acp{peo} and \acp{drpm}~\autocite{SECNAV5400-15D,SECNAV5000-2G}.
  \item[\textbf{Current portfolio authority.}] Since Dec~2025 and the 2026 \ac{don} acquisition restructuring, \acp{pae} are the board-relevant portfolio layer under \ac{asnrdanda}. Public \ac{don} releases describe \acp{pae} as accountable officials for whole portfolios, with authority over program offices and associated technical, contracting, sustainment, and industrial-base functions; about 70 percent of those associated functions and personnel are moving from \acp{syscom} into \acp{pae} as the construct matures~\autocite{DON-PAE-Establishment-2026,DON-PAE-Expansion-May2026}.
  \item[\textbf{Matrix execution.}] The older \ac{peo}/\ac{syscom} matrix remains important for board recall and for live execution during transition: \acp{peo} and \acp{pm} still execute named programs where public pages have not fully caught up; \acp{syscom} (\ac{navsea}, \ac{navwar}, \ac{navair}) still provide technical authority, engineering standards, warfare-center depth, contracting infrastructure, and sustainment support. The PAE construct changes accountability and alignment; it does not give a \ac{pm} permission to ignore \ac{ta} or contracting authority~\autocite{SECNAV5400-15D,DON-PAE-Establishment-2026,DON-PAE-Expansion-May2026}.
  \item[\textbf{Flag/PEO billets (current).}]
        \begin{itemize}
            \item[]
            \item[] \textbf{EDO-led today:} \ac{asnrdanda}'s \ac{pmildep} (\ac{vadm} Seiko Okano), \ac{comnavwar} (\ac{radm} J.\ Kurt Rothenhaus), \ac{peoships} (\ac{rdml} Brian Metcalf), \ac{peocvn} (\ac{radm} Casey Moton), and \ac{peossn} (\ac{radm} Jonathan Rucker) anchor \ac{ta} and acquisition governance for the community~\autocite{navy-flag-roster-2026-05,navy-okano-2025,navy-rothenhaus-2025,navy-metcalf-2025,navy-moton-2025,navy-rucker-2025}.
            \item[] \textbf{Non-EDO leaders (awareness):} \ac{navsea}~08 is currently \ac{admiral} William J.\ Houston (Submarine officer), \ac{peoiws} is \ac{rdml} Thomas J.\ Dickinson (Surface Warfare Acquisition Professional), \ac{peomlb} is \ac{rdml} Kevin P.\ Biehn (Surface Warfare Acquisition Professional), NAVSEA~05 is \ac{rdml} Jay A.\ Seif (Submarine Acquisition Professional), and \ac{peousc} is filled by an \ac{ses} (formerly \ac{rdml} Kevin R.\ Smith)~\autocite{navsea-org-chart-2025}.
        \end{itemize}
        \FlagBilletNote{See Appendix~\ref{app:edo-flags} for the current EDO flag billet roster.}
\end{description}

Figure \ref{fig:chain_of_command_acq_gov} shows the legacy Navy acquisition governance chain of command. Use it to understand the \ac{peo}/\ac{syscom} matrix, then apply the PAE update in Figure~\ref{fig:asnrda_org_chart} and Table~\ref{tab:pae-roster-current}.

\begin{figure}[H]
  \centering\color{SteelDeep}
  \includegraphics[width=0.7\linewidth]{Navy Acquisition Organization.png}
  \caption[Acquisition Governance Chain of Command]{Acquisition Governance Chain of Command.~\srcCite{edo-2-1-1-command-structures-2025}}
  \label{fig:chain_of_command_acq_gov}
\end{figure}

\subsection{ASN (RD\&A) Organization and Portfolio Acquisition Executives}

Current as of 2026-06-04, Mr.\ William F.\ Mahan is performing the duties of \ac{asnrdanda} and serves as the senior acquisition executive for the Navy and Marine Corps. Mr.\ Jason L.\ Potter returned to Principal Civilian Deputy \ac{asnrdanda}, and \ac{pmildep} \ac{vadm} Seiko Okano remains the senior uniformed acquisition deputy supporting \ac{don} governance~\autocite{DON-New-SAE-2026}. Figure~\ref{fig:asnrda_org_chart} captures the board-level PAE organization, Table~\ref{tab:asnrda_leadership} captures top roles, and Table~\ref{tab:pae-roster-current} lists the publicly announced \acp{pae}.

\begin{description}
  \item[\textbf{PAE purpose.}] PAEs align authority with responsibility for outcomes across like portfolios. The reform target is faster delivery, fewer non-value-added reviews, disciplined cost/schedule/performance trades, direct industrial-base management, and cradle-to-grave control across program, technical, contracting, and sustainment functions~\autocite{DON-PAE-Establishment-2026,DON-PAE-Expansion-May2026}.
  \item[\textbf{Board answer.}] Draw \ac{secnav} at the top, \ac{asnrdanda} beneath \ac{secnav}, Principal Civilian Deputy and \ac{pmildep} directly under \ac{asnrdanda}, and the \acp{pae} as the current acquisition portfolio owners. Keep \ac{cno} and \ac{cmc} in the service chain for requirements/readiness and keep \acp{syscom} as the technical, contracting, sustainment, and workforce infrastructure that is being realigned into or in support of the PAEs.
  \item[\textbf{Transition caution.}] The user's PAE organization chart is a useful transition snapshot, but it predates the May~2026 public announcement of PAE Aviation, PAE Mission Systems, and PAE Munitions. Use the current count of nine \acp{pae} unless a newer official chart supersedes it~\autocite{DON-PAE-Expansion-May2026}.
\end{description}

\begin{figure}[H]
  \centering
  \resizebox{\linewidth}{!}{%
  \begin{tikzpicture}
    [
    node distance=1.15cm and 1.15cm,
    top node/.style={rectangle, draw, fill=Vapor, text width=10em, align=center, minimum height=3.5em},
    asn node/.style={rectangle, draw, fill=Gold!40, text width=14em, align=center, minimum height=4.2em},
    deputy node/.style={rectangle, draw, fill=Green!25, text width=11em, align=center, minimum height=3.6em},
    pae node/.style={rectangle, draw, fill=Blue!18, text width=10.5em, align=center, minimum height=3.2em},
    pae group/.style={rectangle, draw=SteelGray, densely dashed, rounded corners, inner sep=0.35cm},
    support node/.style={rectangle, draw, dashed, fill=SteelGray!10, text width=16em, align=center, minimum height=4.0em},
    line/.style={draw, -{Stealth[length=2mm]}, shorten >=1pt, shorten <=1pt},
    support line/.style={draw, dashed, -{Stealth[length=2mm]}, shorten >=1pt, shorten <=1pt}
    ]

    \node[top node] (secnav) {\ac{secnav}};
    \node[top node, left=of secnav, xshift=-2.5cm] (cno) {\ac{cno}};
    \node[top node, right=of secnav, xshift=2.5cm, fill=Red!18] (cmc) {\ac{cmc}};

    \node[asn node, below=of secnav] (asn) {Performing the Duties of\\ \ac{asnrdanda}\\ \textbf{William F.\ Mahan}\\ \small \ac{don} \ac{sae}};
    \node[deputy node, below left=of asn, xshift=-1.6cm] (pcd) {Principal Civilian Deputy\\ \textbf{Jason L.\ Potter}};
    \node[deputy node, below right=of asn, xshift=1.6cm] (pmd) {\ac{pmildep}\\ \textbf{\ac{vadm} Seiko Okano}};

    \node[pae node, below=1.2cm of pcd, xshift=-2.1cm] (aviation) {PAE Aviation};
    \node[pae node, right=of aviation] (missions) {PAE Mission Systems};
    \node[pae node, right=of missions] (munitions) {PAE Munitions};

    \node[pae node, below=of aviation] (industrial) {PAE Industrial Operations};
    \node[pae node, right=of industrial] (maritime) {PAE Maritime};
    \node[pae node, right=of maritime] (undersea) {PAE Undersea\\ / DRPM Submarines};

    \node[pae node, below=of industrial] (ssp) {PAE Strategic\\ Systems Programs};
    \node[pae node, right=of ssp] (ras) {PAE Robotics and\\ Autonomous Systems};
    \node[pae node, right=of ras, fill=Red!18] (marine) {PAE Marine Corps};

    \node[support node, left=of industrial, xshift=-4.0cm] (syscoms) {\acp{syscom}, warfare centers,\\ shipyards, RMCs, NIWCs,\\ contracting, sustainment, and \ac{ta}};
    \begin{scope}[on background layer]
      \node[pae group, fit=(aviation)(missions)(munitions)(industrial)(maritime)(undersea)(ssp)(ras)(marine)] (paeGroup) {};
    \end{scope}
    \coordinate (asnDrop) at ($(asn.south)+(0,-0.55cm)$);
    \coordinate (paeEntry) at ($(paeGroup.north)+(0,0.35cm)$);
    \coordinate (supportEntry) at ($(paeGroup.west)+(-0.35cm,0)$);

    \draw[line] (secnav.south) -- (asn.north);
    \draw[line] (cno.south) |- (asn.west);
    \draw[line] (cmc.south) |- (asn.east);
    \draw[line] (asn.south west) -- (pcd.north);
    \draw[line] (asn.south east) -- (pmd.north);
    \draw[line] (asn.south) -- (asnDrop) -| (paeEntry) -- (paeGroup.north);
    \draw[support line] (syscoms.east) -- ++(0.7cm,0) |- (supportEntry) -- (paeGroup.west);
  \end{tikzpicture}}
  \caption[ASN (RD\&A) and PAE organization transition]{Board-level ASN (RD\&A) and PAE organization transition, updated from the user-provided PAE chart with public announcements through 2026-06-04~\autocite{DON-New-SAE-2026,DON-PAE-Establishment-2026,DON-PAE-Expansion-May2026}.}
  \label{fig:asnrda_org_chart}
\end{figure}

Table~\ref{tab:asnrda_leadership} lists the current top \ac{asnrdanda} roles.

\begin{longtblr}[%
    caption = {ASN (RD\&A) leadership roles},
    label   = {tab:asnrda_leadership},
    entry   = {ASN (RD\&A) leadership},
    remark{\tSource} = {\tabCite{DON-New-SAE-2026}}
  ]{colspec = {@{} llX @{}}}
  \toprule
  {Role}             & {Person}              & {Roles and responsibilities}                                                                                                                                                                                                                      \\
  \midrule
  PTDO ASN (RD\&A) & William F.\ Mahan & Serves as the DON senior acquisition executive for the Navy and Marine Corps while performing the duties of ASN (RD\&A). \\
  Principal Civilian Deputy & Jason L.\ Potter & Provides senior civilian acquisition continuity and returned to this role after serving as PTDO ASN (RD\&A) from July~2025 to May~2026. \\
  PMILDEP            & VADM Seiko Okano      & Provides senior military advice to ASN (RD\&A), helping manage DON research, development, acquisition governance, and the operational perspective across the acquisition enterprise.                                                    \\
  \bottomrule
\end{longtblr}

Table~\ref{tab:pae-roster-current} summarizes the public PAE portfolio map. Use Appendix~\ref{app:pae-construct-mapping} for detailed PAE-to-legacy-organization mapping.

\begin{longtblr}[%
    caption = {Current public PAE roster and board cues},
    label   = {tab:pae-roster-current},
    entry   = {Current public PAE roster},
    remark{\tSource} = {\tabCite{DON-PAE-Establishment-2026,DON-PAE-Expansion-May2026,PAE-Maritime-Homepage-2026,DON-PAE-Mission-Systems-2026,DON-PAE-Munitions-2026,NAVAIR-PAE-Aviation-2026,PAE-Marine-Corps-About-2026,DON-PAE-RAS-MUSV-2026}}
  ]{colspec = {@{} L L X @{}}}
  \toprule
  {PAE}                    & {Public leader / cue}                      & {Board-level scope} \\
  \midrule
  Aviation & VADM John Dougherty & Naval aviation portfolio; public deputy lanes are Carrier Strike, Marine Corps Aviation, and Maritime ISR/NC3. \\
  Industrial Operations & VADM James P.\ Downey & Industrial capacity, naval shipyards, regional maintenance, maintenance contracting, and industrial-base throughput. \\
  Marine Corps & LtGen Eric Austin & Marine Corps ground systems; combines former PEO Land Systems and Marine Corps Systems Command acquisition support functions. \\
  Maritime & Christopher Miller & Conventional surface-ship delivery, including aircraft carriers, combatants, expeditionary and mine warfare, auxiliaries, and modernization/sustainment. \\
  Mission Systems & Jim Day & Central mission integrator for kill-chain closure across PEO C4I, PEO Digital, PEO IWS, PEO MLB, DRPM Overmatch, DRPM LRNFO, DRPM MILE, Minotaur, and mission-system elements across SYSCOMs. \\
  Munitions & Paul Mann & Munitions production and weapons portfolio; public pillars are Weapons Industrial Base, Air Weapons, Surface Weapons, and Advanced Capabilities and Innovation. \\
  Robotics and Autonomous Systems & Public leader not confirmed in this refresh & Hedge and unmanned/autonomous capabilities; current public example is MUSV marketplace at-sea demonstrations. \\
  Strategic Systems Programs & VADM Johnny Wolfe & Sea-based strategic deterrence and strategic weapons; SSP aligned inside the PAE construct. \\
  Undersea / DRPM Submarines & VADM Robert Gaucher & Submarine and undersea delivery; dual-hatted DRPM construct focused on SSN, SSBN, and SSGN production and undersea industrial-base execution. \\
  \bottomrule
\end{longtblr}

\subsection{OPNAV Organization and CNO Principal Assistants}
\begin{description}
  \item[\textbf{OPNAV role.}] \ac{opnav} is the \ac{cno}'s staff within the administrative chain, responsible for translating naval policy, requirements, and resourcing guidance into executable direction across the Fleet and Shore Establishment~\autocite{USC-8032,edo-2-1-1-command-structures-2025}.
\end{description}

\begin{longtblr}[%
    caption = {CNO principal assistants (N-codes)},
    label   = {tab:cno-principal-assistants},
    entry   = {CNO principal assistants},
    remark{\tSource} = {\tabCite{edo-2-1-1-command-structures-2025}}
  ]{colspec = {@{} ll @{}}}
  \toprule
  {N-code} & {Primary responsibilities} \\
  \midrule
  N1       & Manpower, personnel, training, and education policy and readiness. \\
  N2/N6    & Information warfare, intelligence, cyber, command and control, and related information capabilities. \\
  N3/N5    & Operations, plans, and strategy; joint operations and campaign planning. \\
  N4       & Readiness and logistics, including supply, ordnance, energy, and mobility. \\
  N7       & Warfighting development, concepts, and integration. \\
  N8       & Integration of capabilities and resources; PPBE resourcing and assessment. \\
  N9       & Warfare systems requirements and resourcing across major warfighting domains. \\
  \bottomrule
\end{longtblr}

\subsection{Fleet Forces, TYCOMs, and Numbered Fleets}
\begin{description}
  \item[\textbf{USFF.}] \ac{usff} executes the Prepare the Force and Provide the Force functions, generating ready forces for \ac{cno} and \acp{ccdr}, and delivering trained, equipped forces to operational commanders~\autocite{USFF-Mission,edo-2-1-1-command-structures-2025}.
  \item[\textbf{TYCOMs.}] \acp{tycom} supervise specific categories of forces and balance readiness, manning, training, and equipment; they are in the administrative chain of command~\autocite{edo-2-1-1-command-structures-2025}.
  \item[\textbf{Numbered Fleets.}] Numbered Fleet Commanders report to the appropriate \acp{ccdr}; 2nd, 3rd, and 4th Fleets are home-ported in the U.S., 5th, 6th, and 7th Fleets are forward-based, and 10th Fleet operates in cyberspace~\autocite{edo-2-1-1-command-structures-2025}.
\end{description}

\subsection{Echelons and Enterprise Construct}
\begin{description}
  \item[\textbf{Echelons.}] The \ac{cno} is Echelon I; commands reporting to the \ac{cno} are Echelon II; commands reporting to Echelon II are Echelon III, and so on~\autocite{edo-2-1-1-command-structures-2025}.
  \item[\textbf{Enterprise construct.}] Requirements are set by warfighters; providers (\acp{peo}, \acp{syscom}, and warfare centers) design, acquire, build, and maintain; resources are programmed and budgeted to deliver the capability~\autocite{edo-2-1-1-command-structures-2025}.
\end{description}

\subsection{Approving Authorities}

\begin{figure}[H]
  \centering\color{SteelDeep}
  \includegraphics[width=0.7\linewidth]{TA vs PA.png}
  \caption[Technical Authority vs.\ Program Authority]{Technical Authority vs.\ Program Authority alignment.~\srcCite{edo-2-1-5-technical-authority-2025}}
  \label{fig:chain_of_command_auth_roles}
\end{figure}

Figure \ref{fig:chain_of_command_auth_roles} contrasts the roles of \ac{ta} and \ac{pa} in Navy acquisition.
%====================
% End of file
%====================
\IfSubfilesClassLoaded{
  \RenewDocumentCommand{\entryname}{}{\textbf{\color{Modern} Acronym}}
  \RenewDocumentCommand{\descriptionname}{}{\textbf{\color{Modern} Definition}}
  \printnoidxglossary[
    type=\acronymtype,
    title=Acronyms,
    style=long-booktabs
  ]}{}
\subsection{Coursebook cue: EDO flag billets}

Coursebook 2.1.1 asks students to identify the EDO flag billets by name and
responsibility.  Use Appendix~\ref{app:edo-flags} as the live roster for that
board cue because the coursebook roster is a date-stamped training snapshot,
while the appendix is maintained against public Navy flag-officer sources
\autocite{edo-2-1-1-command-structures-2025,navy-flag-roster-2026-04}.

\end{document}
